\documentclass[11pt]{amsart}
% Revised September 22, 2026.
% PDF reproduction: compile to the corresponding public/pdfs path, then run
% python scripts/stamp-pdf-revisions.py in the website checkout (pypdf, reportlab).
% This adds only a first-page margin date; it does not reflow mathematical content.

\usepackage[margin=1.05in]{geometry}
\usepackage{amsmath,amssymb,mathrsfs,mathtools}
\usepackage{enumitem}
\usepackage[dvipsnames]{xcolor}
\usepackage{hyperref}
\usepackage{microtype}

\hypersetup{
  colorlinks=true,
  linkcolor=OliveGreen,
  citecolor=OliveGreen,
  urlcolor=OliveGreen
}

\setlength{\parskip}{0.35em}
\setlength{\parindent}{0pt}

\DeclareMathOperator{\Spec}{Spec}
\DeclareMathOperator{\Proj}{Proj}
\DeclareMathOperator{\Sym}{Sym}
\DeclareMathOperator{\Frac}{Frac}
\DeclareMathOperator{\im}{im}
\DeclareMathOperator{\coker}{coker}
\DeclareMathOperator{\rank}{rank}
\DeclareMathOperator{\Hom}{Hom}
\DeclareMathOperator{\Ann}{Ann}
\DeclareMathOperator{\id}{id}

\newcommand{\OO}{\mathcal O}
\newcommand{\F}{\mathscr F}
\newcommand{\G}{\mathscr G}
\newcommand{\Hh}{\mathscr H}
\newcommand{\A}{\mathscr A}
\newcommand{\E}{\mathscr E}
\newcommand{\I}{\mathscr I}
\newcommand{\U}{\mathscr U}
\newcommand{\V}{\mathscr V}
\newcommand{\tensor}{\otimes}
\newcommand{\into}{\hookrightarrow}
\newcommand{\onto}{\twoheadrightarrow}
\newcommand{\restr}[2]{#1|_{#2}}
\newcommand{\dual}{\vee}

\theoremstyle{plain}
\newtheorem{theorem}{Theorem}[section]
\newtheorem{proposition}[theorem]{Proposition}
\newtheorem{lemma}[theorem]{Lemma}
\newtheorem{corollary}[theorem]{Corollary}

\theoremstyle{definition}
\newtheorem{definition}[theorem]{Definition}
\newtheorem{construction}[theorem]{Construction}

\theoremstyle{remark}
\newtheorem*{strategy}{Strategy}
\newtheorem*{formalization}{Formalization note}
\newtheorem*{remark}{Remark}

\title{Normality, Coherent Extension, Tensor Operations, and Affineness}
\author{Henry Shin}
\date{Notes on Hartshorne, Algebraic Geometry, Chapter II, Section 5, Exercises 14--18}

\begin{document}
\begin{abstract}
These notes reorganize Exercises II.5.14--18 into a single exposition.  The proofs are written with two goals in mind: first, to make the geometric strategy visible, and second, to isolate the algebraic lemmas in a form that can later be ported to Lean.  The common theme is that many global assertions about projective schemes, coherent sheaves, tensor constructions, affine morphisms, and vector bundles become transparent after reducing to affine charts and then proving that the resulting local constructions glue canonically.
\end{abstract}
\maketitle

\section*{Conventions}
All schemes are assumed to be schemes in the sense of Hartshorne. In Section~1, connected schemes are assumed nonempty.  A normal scheme means a scheme whose local rings are integrally closed domains.  If $S$ is a graded ring generated by $S_1$ over $S_0$, we write $S^{(d)} = \bigoplus_{n\geq 0} S_{nd}$ for the $d$-th Veronese subring.  For a closed subscheme $X\subseteq \mathbb P^r_k$, its homogeneous coordinate ring is taken with respect to the given embedding and with the saturated homogeneous ideal defining $X$.

\section{Projective normality after a Veronese embedding}

\begin{strategy}
The core maneuver is to compare the original homogeneous coordinate ring
\[
  S = k[x_0,\ldots,x_r]/I
\]
with the section ring
\[
  S' = \bigoplus_{n\geq 0} \Gamma(X,\OO_X(n)).
\]
For a normal connected projective scheme, $S'$ is not just a convenient enlargement of $S$: it is the normalization of $S$.  Once this is known, the Veronese argument is a degree-bookkeeping argument.  Exercise II.5.9 says that $S$ and $S'$ agree in all sufficiently large degrees; choosing a Veronese degree $d$ beyond that range forces $S^{(d)} = S'^{(d)}$.  Since Veronese subrings of normal graded domains are normal, the $d$-uple embedding is projectively normal.
\end{strategy}

\begin{proposition}[Connected normal noetherian schemes are integral]\label{prop:connected-normal-integral}
Let $X$ be a nonempty connected noetherian scheme such that every local ring $\OO_{X,x}$ is a domain. Then $X$ is integral.
\end{proposition}

\begin{proof}
The condition that all local rings are domains implies that all stalks are reduced, hence $X$ is reduced.  It remains to prove irreducibility.

Since $X$ is noetherian, it has finitely many irreducible components, say
\[
  X = Z_1\cup\cdots\cup Z_m.
\]
If $m>1$, connectedness forces two components to meet.  More explicitly, the graph whose vertices are the $Z_i$ and whose edges record nonempty intersections is connected; otherwise the union of the components in one connected component of the graph would be both closed and open in $X$.  Thus there are $i\neq j$ and a point $x\in Z_i\cap Z_j$.

Choose an affine open neighbourhood $U=\Spec A$ of $x$, and let $\mathfrak q\subset A$ be the prime corresponding to $x$.  The nonempty closed subsets $Z_i\cap U$ and $Z_j\cap U$ contain irreducible components of $U$, corresponding to distinct minimal primes $\mathfrak p_i,\mathfrak p_j\subset A$ contained in $\mathfrak q$.  After localizing at $\mathfrak q$, the ring $A_{\mathfrak q}=\OO_{X,x}$ has at least the two distinct minimal primes $\mathfrak p_iA_{\mathfrak q}$ and $\mathfrak p_jA_{\mathfrak q}$.  This contradicts the assumption that $\OO_{X,x}$ is a domain.  Hence $m=1$, so $X$ is irreducible.  Since $X$ is also reduced, it is integral.
\end{proof}

\begin{proposition}[The section ring is the normalization]\label{prop:section-ring-normalization}
Let $k$ be an algebraically closed field, and let $X\subseteq \mathbb P^r_k$ be a connected normal closed subscheme.  Let
\[
  S = k[x_0,\ldots,x_r]/I
\]
be its homogeneous coordinate ring, and set
\[
  S' = \bigoplus_{n\geq 0} \Gamma(X,\OO_X(n)).
\]
Then $S$ is a domain and $S'$ is the integral closure of $S$ in $\Frac(S)$.
\end{proposition}

\begin{proof}
By Proposition \ref{prop:connected-normal-integral}, $X$ is integral.  Since $X$ is an integral closed subscheme of projective space, its saturated homogeneous ideal $I$ is prime.  Indeed, if $f,g$ are homogeneous forms and $fg$ vanishes on $X$, then
\[
  X \subseteq V(fg) = V(f)\cup V(g).
\]
Irreducibility gives $X\subseteq V(f)$ or $X\subseteq V(g)$, and reducedness then gives $f\in I$ or $g\in I$.  For a homogeneous ideal this homogeneous-element criterion is equivalent to primeness.  Thus $S$ is a domain.

Let $K=\Frac(S)$.  For each $i$ with $D_+(x_i)\neq\varnothing$, we have
\[
  D_+(x_i) \cong \Spec S_{(x_i)}.
\]
For every $n\geq 0$,
\[
  \Gamma(D_+(x_i),\OO_X(n)) = (S_{x_i})_n = x_i^nS_{(x_i)}
\]
inside the homogeneous fraction field of $S$.  Since the $D_+(x_i)$ cover $X$, a global section is exactly a collection of local sections agreeing on overlaps; under the above embeddings this gives
\[
  \Gamma(X,\OO_X(n))
  = \bigcap_i x_i^n S_{(x_i)}
\]
inside $K$.  Hence
\[
  S' = \bigcap_i \bigoplus_{n\geq 0} x_i^nS_{(x_i)}.
\]

Now $S_{(x_i)}$ is normal.  Its localizations at prime ideals are precisely the local rings of $X$ at points of $D_+(x_i)$, and normality is local.  The graded ring
\[
  B_i := \bigoplus_{n\geq 0}x_i^nS_{(x_i)}
\]
identifies with the polynomial ring $S_{(x_i)}[T]$ by sending $T$ to $x_i$; the uniqueness of homogeneous decomposition in $S_{x_i}$ gives injectivity.  Since a polynomial ring over a normal domain is normal, each $B_i$ is normal.  An intersection of normal domains inside a common fraction field is normal, so $S'$ is normal.

To prove that $S\to S'$ is integral, apply Exercise II.5.9 directly to the finitely generated graded module $S$: there is an $N\geq 1$ such that $S_n=S'_n$ for all $n\geq N$. Each $S'_n$ is finite-dimensional over $k$, by finiteness of global sections of a coherent sheaf on a projective scheme. Choose a $k$-basis of each $S'_n$ with $0\leq n<N$. These finitely many elements, together with $1$, generate $S'$ as an $S$-module: the low-degree pieces are spanned by the chosen bases, and every higher-degree piece already lies in $S$. Thus $S'$ is finite over $S$, hence integral over $S$. These inputs do not require normality of the section ring or a Veronese embedding. Therefore $S'\subseteq \overline S$, where $\overline S$ denotes the integral closure of $S$ in $K$.

Conversely, if $u\in K$ is integral over $S$, then $u$ is also integral over $S'$ because $S\subseteq S'$.  Since $S'$ is integrally closed and has fraction field $K$, we get $u\in S'$.  Hence $\overline S\subseteq S'$, proving $S'=\overline S$.
\end{proof}

\begin{proposition}[Agreement in high degree]\label{prop:high-degree-agreement}
With $S$ and $S'$ as above, there exists $N$ such that
\[
  S_n = S'_n
\]
for every $n\geq N$.
\end{proposition}

\begin{proof}
Apply Exercise II.5.9 to the finitely generated graded $S$-module $M=S$.  Since $\widetilde S\cong \OO_X$, the natural map
\[
  S_n \longrightarrow \Gamma(X,\OO_X(n)) = S'_n
\]
is an isomorphism for all sufficiently large $n$.
\end{proof}

\begin{lemma}[Normality of Veronese subrings]\label{lem:veronese-normal}
Let $R=\bigoplus_{n\geq 0}R_n$ be a normal $\mathbb N$-graded domain.  Then $R^{(d)}=\bigoplus_{n\geq 0}R_{nd}$ is normal for every $d>0$.
\end{lemma}

\begin{proof}
Let $L=\Frac(R^{(d)})$.  Suppose $u\in L$ is integral over $R^{(d)}$.  Since $L\subseteq \Frac(R)$ and $R^{(d)}\subseteq R$, the same monic equation shows that $u$ is integral over $R$.  Normality of $R$ gives $u\in R$.

It remains to prove that $u$ has only homogeneous components of degree divisible by $d$.  Write
\[
  u = u_0+u_1+\cdots+u_m,
\]
where $u_j\in R_j$.  Since $u\in \Frac(R^{(d)})$, write $u=a/b$ with $a,b\in R^{(d)}$ and $b\neq 0$.  Let $b_\ell$ be the lowest-degree nonzero homogeneous component of $b$; its degree $\ell$ is divisible by $d$.

Assume for contradiction that some $u_j$ has $j$ not divisible by $d$, and choose the least such $j$.  In the equality $a=ub$, the component of degree $j+\ell$ is exactly $u_jb_\ell$.  Indeed, any other contribution would be $u_{j'}b_{\ell'}$ with $j'<j$ and $\ell'>\ell$; by minimality of $j$, such $j'$ is divisible by $d$, and $\ell'$ is divisible by $d$, so $j'+\ell'$ is divisible by $d$, while $j+\ell$ is not.  But $a\in R^{(d)}$, so its component of degree $j+\ell$ is zero.  Hence $u_jb_\ell=0$, contradicting that $R$ is a domain.  Therefore all nonzero $u_j$ have $d\mid j$, and $u\in R^{(d)}$.
\end{proof}

\begin{theorem}[Projective normality after a Veronese embedding]\label{thm:veronese-projectively-normal}
Let $k$ be algebraically closed and let $X\subseteq \mathbb P^r_k$ be connected and normal.  Then for all sufficiently large $d>0$, the $d$-uple embedding of $X$ is projectively normal.
\end{theorem}

\begin{proof}
By Proposition \ref{prop:high-degree-agreement}, choose $N$ such that $S_n=S'_n$ for all $n\geq N$.  Also $S'_0=k$: by Proposition \ref{prop:section-ring-normalization}, $S'$ is integral over $S$.  If $a\in S'_0$ satisfies a monic integral equation over $S$, taking the degree-zero part of that equation gives a monic algebraic equation for $a$ over $S_0=k$.  Since $k$ is algebraically closed, $a\in k$, so $S'_0=k$.

Choose $d\geq N$.  Then for every $m\geq 1$,
\[
  S_{md} = S'_{md},
\]
and in degree $0$ both rings are $k$.  Thus
\[
  S^{(d)} = S'^{(d)}.
\]
The ring $S'$ is normal by Proposition \ref{prop:section-ring-normalization}, so $S'^{(d)}$ is normal by Lemma \ref{lem:veronese-normal}.  Hence $S^{(d)}$ is an integrally closed domain.  Since $S^{(d)}$ is the homogeneous coordinate ring of the image of $X$ under the $d$-uple embedding, that embedding is projectively normal.
\end{proof}

\begin{corollary}[Intrinsic criterion for projective normality]\label{cor:projective-normality-criterion}
Let $X\subseteq \mathbb P^r_k$ be a connected closed subscheme over an algebraically closed field, with homogeneous coordinate ring $S$.  Then $X$ is projectively normal for the given embedding if and only if $X$ is normal and, for every $n\geq 0$, the natural restriction map
\[
  \Gamma(\mathbb P^r_k,\OO_{\mathbb P^r}(n))
  \longrightarrow
  \Gamma(X,\OO_X(n))
\]
is surjective.
\end{corollary}

\begin{proof}
Assume first that $X$ is projectively normal.  Then $S$ is an integrally closed domain.  Localizing an integrally closed domain preserves integral closedness, and taking the degree-zero part of a homogeneous localization preserves normality: if an element of the fraction field of the degree-zero part is integral over the degree-zero part, it is integral over the whole localization and hence lies in the localization; the grading then forces it to have degree zero.  Therefore the local rings of $X=\Proj S$ are normal domains, so $X$ is normal.

By Proposition \ref{prop:section-ring-normalization}, applied to this normal $X$, the section ring $S'=\bigoplus_n\Gamma(X,\OO_X(n))$ is the integral closure of $S$.  Since $S$ is already integrally closed, $S=S'$.  The map
\[
  k[x_0,\ldots,x_r]_n \longrightarrow S_n = S'_n = \Gamma(X,\OO_X(n))
\]
is therefore surjective for every $n$.

Conversely, assume that $X$ is normal and that all the restriction maps are surjective.  The map from the polynomial ring to $S'$ factors through the homogeneous coordinate ring $S$.  Surjectivity of the restriction maps gives a surjection $S\onto S'$.  This map is also injective because $S$ is defined using the saturated homogeneous ideal of $X$: a homogeneous element of $S$ mapping to the zero global section of $\OO_X(n)$ vanishes on each $D_+(x_i)$, hence belongs to the saturation, and is therefore zero in $S$.  Thus $S\cong S'$.  By Proposition \ref{prop:section-ring-normalization}, $S'$ is normal, so $S$ is an integrally closed domain.  Hence $X$ is projectively normal.
\end{proof}

\begin{formalization}
For Lean, the safest route is to split this section into lemmas: connected normal noetherian implies irreducible; the section ring embeds as an intersection of the affine chart rings; intersections of normal subrings in a common fraction field are normal; and Veronese subrings of normal graded domains are normal.  The Veronese lemma above is especially well suited to formalization because it is a finite homogeneous-degree argument.
\end{formalization}

\section{Extension of coherent sheaves}

\begin{strategy}
The organizing principle is that noetherian hypotheses let us pass from arbitrary quasi-coherent objects to coherent approximations.  On an affine noetherian scheme, a quasi-coherent sheaf is the union of the coherent sheaves coming from finitely generated submodules.  Once that is known, extension from an open subset is obtained by pushing forward to the ambient affine scheme and then choosing one coherent subsheaf large enough to contain finitely many local generators.  The global noetherian case is obtained by extending over one affine chart at a time.
\end{strategy}

\begin{proposition}[Affine union lemma]\label{prop:affine-union}
Let $X=\Spec A$ be noetherian affine.  Every quasi-coherent sheaf on $X$ is the union of its coherent subsheaves.
\end{proposition}

\begin{proof}
Let $\F=\widetilde M$ for an $A$-module $M$.  Write $M=\bigcup_\lambda M_\lambda$, where $M_\lambda$ runs through the finitely generated $A$-submodules of $M$.  The index set is directed under sum.  Since $A$ is noetherian, each $\widetilde{M_\lambda}$ is coherent, and each is naturally a subsheaf of $\widetilde M$.

The sectionwise directed union of these subsheaves is a sheaf: every open subset of a noetherian space is quasi-compact, so any gluing problem reduces to a finite subcover, and directedness puts the finitely many sections in one subsheaf. On stalks we have
\[
  \left(\bigcup_\lambda \widetilde{M_\lambda}\right)_\mathfrak p
  = \bigcup_\lambda (M_\lambda)_\mathfrak p
  = M_\mathfrak p
  = \F_\mathfrak p.
\]
Thus the union is equal to $\F$.
\end{proof}

\begin{proposition}[Extending coherent sheaves from an open subset of an affine]\label{prop:affine-extension}
Let $X=\Spec A$ be noetherian affine, let $U\subseteq X$ be open, and let $\F$ be coherent on $U$.  Then there is a coherent sheaf $\F'$ on $X$ such that $\F'|_U\cong \F$.
\end{proposition}

\begin{proof}
Let $i:U\into X$ be the inclusion.  Since $X$ is noetherian, $U$ is quasi-compact, and $i_*\F$ is quasi-coherent on $X$.  By Proposition \ref{prop:affine-union},
\[
  i_*\F = \bigcup_\lambda \Hh_\lambda
\]
with each $\Hh_\lambda$ coherent on $X$.

Choose a finite affine open cover $U=\bigcup_{j=1}^m U_j$ such that $\F|_{U_j}$ is generated by finitely many sections $s_{j1},\ldots,s_{jN_j}$.  Each $s_{ja}$ is a section of $i_*\F$ over $U_j$, hence belongs to $\Hh_{\lambda(j,a)}(U_j)$ for some index $\lambda(j,a)$. Indeed, stalkwise membership gives a cover of $U_j$ on which this holds; quasi-compactness reduces it to a finite cover, and directedness puts all these local sections in one $\Hh_\lambda$, where they glue.  Since the family is directed and there are only finitely many such sections, choose one index $\lambda_0$ dominating all of them.

Then $\Hh_{\lambda_0}|_U\subseteq \F$ and, on each $U_j$, its image contains a generating set for $\F|_{U_j}$.  Hence $\Hh_{\lambda_0}|_U=\F$.  Taking $\F'=\Hh_{\lambda_0}$ proves the claim.
\end{proof}

\begin{proposition}[Affine extension inside a quasi-coherent sheaf]\label{prop:affine-extension-sub}
Let $X=\Spec A$ be noetherian affine, let $U\subseteq X$ be open, let $\F$ be coherent on $U$, and let $\G$ be quasi-coherent on $X$ with an inclusion
\[
  \F \subseteq \G|_U.
\]
Then there is a coherent subsheaf $\F'\subseteq \G$ such that $\F'|_U=\F$.
\end{proposition}

\begin{proof}
Let $i:U\into X$ be the inclusion.  There is a natural restriction morphism
\[
  \rho:\G\longrightarrow i_*(\G|_U).
\]
Define
\[
  \Hh := \rho^{-1}(i_*\F)\subseteq \G.
\]
Equivalently, $\Hh$ is the kernel of the composite
\[
  \G \longrightarrow i_*(\G|_U) \longrightarrow i_*((\G|_U)/\F).
\]
Because $i$ is a quasi-compact open immersion and $X$ is affine noetherian, the sheaves appearing here are quasi-coherent; hence $\Hh$ is quasi-coherent.  Moreover $\Hh|_U=\F$.

Now repeat the finite-generator argument from Proposition \ref{prop:affine-extension}.  Write $\Hh$ as a union of coherent subsheaves of $\G$ using Proposition \ref{prop:affine-union}; choose one coherent subsheaf large enough to contain local generators for $\F$ on a finite affine cover of $U$.  Its restriction to $U$ is exactly $\F$.
\end{proof}

\begin{theorem}[Global extension inside a quasi-coherent sheaf]\label{thm:global-extension-sub}
Let $X$ be noetherian, let $U\subseteq X$ be open, let $\F$ be coherent on $U$, and let $\G$ be quasi-coherent on $X$ with $\F\subseteq \G|_U$.  Then there is a coherent subsheaf $\F'\subseteq \G$ such that $\F'|_U=\F$.
\end{theorem}

\begin{proof}
Choose a finite affine open cover $X=V_1\cup\cdots\cup V_m$.  We extend over the $V_i$ one at a time.

Set $W_0=U$ and $\F_0=\F$.  Suppose by induction that $\F_j\subseteq \G|_{W_j}$ is coherent, where
\[
  W_j = U\cup V_1\cup\cdots\cup V_j,
\]
and that $\F_j|_U=\F$.  Consider the affine scheme $V_{j+1}$ and the open subset
\[
  W_j\cap V_{j+1}\subseteq V_{j+1}.
\]
The sheaf $\F_j|_{W_j\cap V_{j+1}}$ is coherent and is a subsheaf of $\G|_{W_j\cap V_{j+1}}$.  By Proposition \ref{prop:affine-extension-sub}, it extends to a coherent subsheaf
\[
  \Hh_{j+1}\subseteq \G|_{V_{j+1}}
\]
whose restriction to $W_j\cap V_{j+1}$ equals $\F_j|_{W_j\cap V_{j+1}}$.

Thus $\F_j$ on $W_j$ and $\Hh_{j+1}$ on $V_{j+1}$ glue along their common restriction to give a coherent subsheaf
\[
  \F_{j+1}\subseteq \G|_{W_j\cup V_{j+1}}.
\]
After finitely many steps we get the desired coherent subsheaf on all of $X$.
\end{proof}

\begin{corollary}[Extension of coherent sheaves]\label{cor:extension-coherent}
Let $X$ be noetherian, let $U\subseteq X$ be open, and let $\F$ be coherent on $U$.  Then there is a coherent sheaf $\F'$ on $X$ with $\F'|_U\cong \F$.
\end{corollary}

\begin{proof}
Take $\G=i_*\F$, where $i:U\into X$ is the inclusion.  Since $X$ is noetherian, $i$ is quasi-compact, and $i_*\F$ is quasi-coherent.  Apply Theorem \ref{thm:global-extension-sub}.
\end{proof}

\begin{corollary}[Global union of coherent subsheaves]\label{cor:global-union}
On a noetherian scheme, every quasi-coherent sheaf is the union of its coherent subsheaves.
\end{corollary}

\begin{proof}
Let $\G$ be quasi-coherent on $X$.  It is enough to show that every local section of $\G$ lies in some coherent subsheaf of $\G$.

Let $s\in \Gamma(U,\G)$ for an open $U\subseteq X$.  The image of the map
\[
  \OO_U \longrightarrow \G|_U,\qquad 1\longmapsto s
\]
is a coherent subsheaf of $\G|_U$: on affine opens it is the cyclic module generated by $s$, and the kernel is an ideal in a noetherian ring, hence finitely generated.  By Theorem \ref{thm:global-extension-sub}, this cyclic coherent subsheaf extends to a coherent subsheaf of $\G$ on $X$.  Therefore every local section is contained in a coherent subsheaf, which is exactly the assertion that $\G$ is the union of its coherent subsheaves.
\end{proof}

\begin{formalization}
The clean formal pattern is: prove the affine module statement first, then use directed unions and finite local generation.  The inverse-image sheaf $\rho^{-1}(i_*\F)$ in Proposition \ref{prop:affine-extension-sub} should be implemented as the kernel of $\G\to i_*((\G|_U)/\F)$; this avoids ambiguity about preimages of subsheaves.
\end{formalization}

\section{Tensor operations on sheaves}

\begin{strategy}
Most tensor-operation statements are local.  For locally free sheaves, choose a basis and reduce to the corresponding calculation for a free module.  For filtrations coming from a short exact sequence, the important point is to define the filtration canonically, not by choosing a splitting.  Local splittings are then used only to identify the associated graded pieces.
\end{strategy}

\begin{proposition}[Ranks of tensor, symmetric, and exterior powers]\label{prop:tensor-ranks}
Let $\F$ be locally free of rank $n\geq 1$ on a ringed space $X$. For $r\geq 0$, the sheaves $T^r(\F)$, $\Sym^r(\F)$, and $\bigwedge^r\F$ are locally free of ranks
\[
  n^r,\qquad \binom{n+r-1}{r},\qquad \binom nr
\]
respectively, with $\binom nr=0$ for $r>n$. If $n=0$, all three degree-zero powers are $\OO_X$, and every positive-degree power is zero.
\end{proposition}

\begin{proof}
The assertion is local on $X$.  On an open set where $\F\cong \OO_X^{\oplus n}$ with basis $e_1,\ldots,e_n$, the tensor power $T^r(\F)$ has basis
\[
  e_{i_1}\otimes\cdots\otimes e_{i_r},\qquad 1\leq i_j\leq n,
\]
so its rank is $n^r$.  The symmetric power has basis the monomials
\[
  e_1^{a_1}\cdots e_n^{a_n},\qquad a_1+\cdots+a_n=r,
\]
so its rank is $\binom{n+r-1}{r}$.  The exterior power has basis
\[
  e_{i_1}\wedge\cdots\wedge e_{i_r},\qquad 1\leq i_1<\cdots<i_r\leq n,
\]
so its rank is $\binom nr$.
\end{proof}

\begin{proposition}[Perfect exterior pairing]\label{prop:exterior-perfect-pairing}
Let $\F$ be locally free of rank $n$.  For every $0\leq r\leq n$, the wedge product
\[
  \bigwedge^r\F \otimes \bigwedge^{n-r}\F
  \longrightarrow
  \bigwedge^n\F
\]
is a perfect pairing.  Equivalently, it induces an isomorphism
\[
  \bigwedge^r\F
  \cong
  \left(\bigwedge^{n-r}\F\right)^\dual\otimes \bigwedge^n\F.
\]
In particular, if $\rank \F=2$, then
\[
  \F \cong \F^\dual\otimes \bigwedge^2\F.
\]
\end{proposition}

\begin{proof}
The claim is local.  On a trivializing open set with basis $e_1,\ldots,e_n$, the basis element
\[
  e_I=e_{i_1}\wedge\cdots\wedge e_{i_r}
\]
of $\bigwedge^r\F$ pairs nontrivially only with the complementary basis element
\[
  e_{I^c}=e_{j_1}\wedge\cdots\wedge e_{j_{n-r}}
\]
of $\bigwedge^{n-r}\F$, and the product is $\pm e_1\wedge\cdots\wedge e_n$.  Thus the matrix of the pairing in these bases is a signed permutation matrix, hence invertible.  The local isomorphisms agree on overlaps because the wedge product is canonical.
\end{proof}

\begin{theorem}[Filtration of symmetric powers]\label{thm:symmetric-filtration}
Let
\[
  0\longrightarrow \F'\longrightarrow \F\longrightarrow \F''\longrightarrow 0
\]
be an exact sequence of locally free sheaves.  For every $r\geq 0$, there is a finite decreasing filtration
\[
  \Sym^r\F = F^0 \supseteq F^1 \supseteq \cdots \supseteq F^{r}\supseteq F^{r+1}=0
\]
whose associated graded pieces are canonically isomorphic to
\[
  F^p/F^{p+1}
  \cong
  \Sym^p\F'\otimes \Sym^{r-p}\F''
\]
for $0\leq p\leq r$.
\end{theorem}

\begin{proof}
Define $F^p$ canonically as the image of the multiplication morphism
\[
  \Sym^p\F'\otimes \Sym^{r-p}\F
  \longrightarrow
  \Sym^r\F.
\]
This gives a decreasing filtration because the image obtained using at least $p+1$ factors from $\F'$ is contained in the image obtained using at least $p$ factors from $\F'$.

To identify the quotients, work locally on an open set where the short exact sequence splits:
\[
  \F \cong \F'\oplus \F''.
\]
Then the standard decomposition of symmetric powers gives
\[
  \Sym^r(\F'\oplus\F'')
  \cong
  \bigoplus_{i=0}^r \Sym^i\F'\otimes \Sym^{r-i}\F''.
\]
Under this decomposition, $F^p$ is the direct sum of the terms with $i\geq p$:
\[
  F^p|_U
  \cong
  \bigoplus_{i=p}^r \Sym^i\F'|_U\otimes \Sym^{r-i}\F''|_U.
\]
Therefore
\[
  F^p/F^{p+1}
  \cong
  \Sym^p\F'\otimes \Sym^{r-p}\F''
\]
locally. The quotient map sends a product of $p$ sections of $\F'$ and $r-p$ local lifts of sections of $\F''$ to its class in $F^p/F^{p+1}$. Replacing any lift changes the product by a term with at least $p+1$ factors in $\F'$, hence by zero in the quotient. Thus this map is independent of the lifts and of the splitting. The local isomorphisms therefore agree on overlaps and glue.
\end{proof}

\begin{theorem}[Filtration of exterior powers]\label{thm:exterior-filtration}
Let
\[
  0\longrightarrow \F'\longrightarrow \F\longrightarrow \F''\longrightarrow 0
\]
be an exact sequence of locally free sheaves.  For every $r\geq 0$, there is a finite decreasing filtration
\[
  \bigwedge^r\F = E^0 \supseteq E^1 \supseteq \cdots \supseteq E^r\supseteq E^{r+1}=0
\]
with associated graded pieces
\[
  E^p/E^{p+1}
  \cong
  \bigwedge^p\F'\otimes \bigwedge^{r-p}\F''.
\]
In particular, if $\F'$, $\F$, and $\F''$ have ranks $n'$, $n$, and $n''$, with $n=n'+n''$, then
\[
  \bigwedge^n\F
  \cong
  \bigwedge^{n'}\F'\otimes \bigwedge^{n''}\F''.
\]
\end{theorem}

\begin{proof}
Define $E^p$ to be the image of the canonical wedge multiplication map
\[
  \bigwedge^p\F'\otimes \bigwedge^{r-p}\F
  \longrightarrow
  \bigwedge^r\F.
\]
The same argument as in Theorem \ref{thm:symmetric-filtration} applies.  Locally split $\F\cong \F'\oplus\F''$ and use
\[
  \bigwedge^r(\F'\oplus\F'')
  \cong
  \bigoplus_{i=0}^r \bigwedge^i\F'\otimes \bigwedge^{r-i}\F''.
\]
Then $E^p$ is locally the direct sum of the terms with $i\geq p$, and the quotient $E^p/E^{p+1}$ is the term with $i=p$.  Taking $r=n$ leaves only the top-rank term.
\end{proof}

\begin{proposition}[Pullback commutes with tensor operations]\label{prop:pullback-tensor-ops}
Let $f:X\to Y$ be a morphism of ringed spaces, and let $\F$ be an $\OO_Y$-module.  Then there are natural isomorphisms
\[
  f^*(T^n\F)\cong T^n(f^*\F),\qquad
  f^*(\Sym^n\F)\cong \Sym^n(f^*\F),\qquad
  f^*\left(\bigwedge^n\F\right)\cong \bigwedge^n(f^*\F).
\]
The same holds for the full tensor, symmetric, and exterior algebras.
\end{proposition}

\begin{proof}
Recall that
\[
  f^*\F = f^{-1}\F\otimes_{f^{-1}\OO_Y}\OO_X.
\]
The inverse image functor commutes with tensor products, and extension of scalars commutes with tensor products.  Therefore
\[
  f^*(\F^{\otimes n})
  \cong
  (f^*\F)^{\otimes n}.
\]
This proves the statement for tensor powers and hence for the tensor algebra.

For symmetric and exterior powers, use the quotient descriptions.  The sheaf $\Sym^n\F$ is the quotient of $T^n\F$ by the subsheaf generated locally by tensors imposing commutativity, and $\bigwedge^n\F$ is the quotient by the subsheaf generated locally by tensors imposing alternation.  Pullback is right exact, and under the tensor-power isomorphism just constructed, the images of the pulled-back relation maps are exactly the corresponding relation subsheaves for $f^*\F$.  Hence the induced quotients are canonically isomorphic.
\end{proof}

\begin{formalization}
The filtration proof above is a streamlined version of the local-splitting idea: define the filtration globally as an image, then use local splittings only to compute the quotient.  The product modulo $F^{p+1}$ is independent of the chosen lifts; local splittings then show that this canonical map is an isomorphism.
\end{formalization}

\section{Affine morphisms and relative Spec}

\begin{strategy}
Affine morphisms are controlled by affine charts on the target.  The main point in part (a) is that if preimages of one affine cover are affine, then preimages of all affine opens are affine.  This uses the affine criterion: a scheme covered by finitely many principal opens coming from global functions is affine if those principal opens are affine.  Relative $\Spec$ is then the global object obtained by gluing the affine schemes $\Spec \A(V)$ as $V$ ranges over affine opens of the base.
\end{strategy}

\begin{theorem}[Affine morphisms can be tested on every affine open]\label{thm:affine-test}
A morphism $f:X\to Y$ is affine if and only if $f^{-1}(V)$ is affine for every open affine $V\subseteq Y$.
\end{theorem}

\begin{proof}
One direction is immediate from the definition.  Conversely, suppose $Y$ has an affine open cover $\{Y_i\}$ such that $f^{-1}(Y_i)$ is affine for all $i$.  Let $V\subseteq Y$ be open affine.  We prove that $f^{-1}(V)$ is affine.

Cover each intersection $V\cap Y_i$ by open subsets that are principal affine opens in both $V$ and $Y_i$.  Since $V$ is quasi-compact, choose finitely many such opens
\[
  V_{1},\ldots,V_m
\]
covering $V$.  Write $V_j=D(g_j)$ inside $V$, with $g_j\in \Gamma(V,\OO_Y)$.  The $g_j$ generate the unit ideal because the $D(g_j)$ cover $V$.

Let $\overline g_j=f^\#(g_j)\in \Gamma(f^{-1}(V),\OO_X)$.  Then
\[
  f^{-1}(V_j) = (f^{-1}(V))_{\overline g_j}.
\]
Each $V_j$ is contained in some $Y_i$ and is principal affine in that $Y_i$, so $f^{-1}(V_j)$ is a principal open subset of the affine scheme $f^{-1}(Y_i)$; hence it is affine.  The functions $\overline g_j$ generate the unit ideal on $f^{-1}(V)$, and the principal opens $(f^{-1}(V))_{\overline g_j}$ are affine.  By the affine criterion of Exercise II.2.17, $f^{-1}(V)$ is affine.
\end{proof}

\begin{proposition}[Basic properties of affine morphisms]\label{prop:affine-basic-properties}
Affine morphisms are quasi-compact and separated.  Every finite morphism is affine.
\end{proposition}

\begin{proof}
If $f:X\to Y$ is affine and $V\subseteq Y$ is affine, then $f^{-1}(V)$ is affine by Theorem \ref{thm:affine-test}, hence quasi-compact.  Since quasi-compactness can be checked on an affine open cover of the target, $f$ is quasi-compact.

Separatedness is local on the target.  Thus assume $Y=\Spec A$ and $X=\Spec B$.  Then
\[
  X\times_Y X \cong \Spec(B\otimes_A B),
\]
and the diagonal morphism corresponds to the multiplication map
\[
  B\otimes_A B \longrightarrow B,
  \qquad b\otimes b'\longmapsto bb'.
\]
This map is surjective, so the diagonal is a closed immersion.  Hence $f$ is separated.

Finally, if $f$ is finite, then by definition the inverse image of each affine open $V=\Spec A\subseteq Y$ is affine of the form $\Spec B$, where $B$ is a finite $A$-algebra.  Hence $f$ is affine.
\end{proof}

\begin{construction}[Relative Spec]\label{cons:relative-spec}
Let $Y$ be a scheme and let $\A$ be a quasi-coherent sheaf of $\OO_Y$-algebras.  There is a scheme
\[
  X=\Spec_Y\A
\]
and a morphism $f:X\to Y$ characterized by the following property: for every open affine $V\subseteq Y$,
\[
  f^{-1}(V)\cong \Spec \A(V),
\]
and for every inclusion of open affines $U\subseteq V$, the morphism
\[
  f^{-1}(U)\into f^{-1}(V)
\]
corresponds to the restriction homomorphism $\A(V)\to \A(U)$.
\end{construction}

\begin{proof}
For each open affine $V=\Spec R\subseteq Y$, the algebra $\A(V)$ is an $R$-algebra, so there is a natural morphism
\[
  \Spec\A(V)\longrightarrow \Spec R=V.
\]
If $U=D(g)\subseteq V$ is principal affine, quasi-coherence gives
\[
  \A(U)\cong \A(V)_g.
\]
Thus $\Spec\A(U)$ is the principal open subset of $\Spec\A(V)$ where the image of $g$ is nonzero.

For two affine opens $V,W\subseteq Y$, cover $V\cap W$ by open subsets that are principal affine in both $V$ and $W$.  The identifications above identify the corresponding principal opens in $\Spec\A(V)$ and $\Spec\A(W)$ canonically.  These identifications agree on triple overlaps because they are induced by restriction maps of the sheaf of algebras $\A$.  Therefore the schemes $\Spec\A(V)$ glue to a scheme $X$, and the maps to $V\subseteq Y$ glue to a morphism $f:X\to Y$.

The stated compatibility with inclusions of affine opens is built into the gluing.  Uniqueness follows because the affine opens $f^{-1}(V)$ cover $X$ and all transition maps are forced by the restriction homomorphisms of $\A$.
\end{proof}

\begin{proposition}[Structure sheaf of relative Spec]\label{prop:relative-spec-structure}
Let $f:X=\Spec_Y\A\to Y$ be the relative spectrum of a quasi-coherent $\OO_Y$-algebra.  Then $f$ is affine and
\[
  f_*\OO_X\cong \A
\]
as sheaves of $\OO_Y$-algebras.  Conversely, if $f:X\to Y$ is affine and $\A=f_*\OO_X$, then $\A$ is a quasi-coherent $\OO_Y$-algebra and
\[
  X\cong \Spec_Y\A
\]
over $Y$.
\end{proposition}

\begin{proof}
For every affine open $V\subseteq Y$, the construction gives
\[
  f^{-1}(V)=\Spec\A(V),
\]
so $f$ is affine.  On such a $V$,
\[
  (f_*\OO_X)(V)=\Gamma(f^{-1}(V),\OO_X)=\A(V).
\]
These identifications are compatible with restrictions, hence glue to $f_*\OO_X\cong \A$.

Conversely, suppose $f$ is affine. Write $V=\Spec R\subseteq Y$ and $f^{-1}(V)=\Spec B$, and set $\A=f_*\OO_X$. For every $g\in R$, the inverse image of $D(g)$ is $\Spec B_g$, so
\[
  \A(D(g))=\Gamma(\Spec B_g,\OO_X)=B_g.
\]
These identifications commute with restriction. Hence $\A|_V\cong\widetilde B$, proving quasi-coherence directly from the affine localization criterion. Relative $\Spec$ can therefore be applied to $\A$. On each affine $V$ we have
\[
  f^{-1}(V)
  \cong
  \Spec \Gamma(f^{-1}(V),\OO_X)
  =
  \Spec (f_*\OO_X)(V).
\]
The restriction maps are exactly the maps on coordinate rings induced by open immersions.  Therefore $X$ satisfies the defining property of $\Spec_Y(f_*\OO_X)$.
\end{proof}

\begin{theorem}[Quasi-coherent sheaves on an affine morphism]\label{thm:affine-qc-equivalence}
Let $f:X\to Y$ be affine and let $\A=f_*\OO_X$.  Then $f_*$ induces an equivalence between quasi-coherent $\OO_X$-modules and quasi-coherent $\A$-modules on $Y$.
\end{theorem}

\begin{proof}
Work locally on affine opens $V=\Spec R\subseteq Y$.  Then $f^{-1}(V)=\Spec B$, where $B=\A(V)$.  Quasi-coherent sheaves on $f^{-1}(V)$ are equivalent to $B$-modules, and quasi-coherent $\A|_V$-modules are also equivalent to $B$-modules.  Under these equivalences, $f_*$ is the identity on the underlying $B$-module.

Given a quasi-coherent $\A$-module $\mathscr M$ on $Y$, define a quasi-coherent sheaf $\widetilde{\mathscr M}$ on $X$ by taking, on each $f^{-1}(V)=\Spec\A(V)$, the sheaf associated to the $\A(V)$-module $\mathscr M(V)$.  The quasi-coherence of $\mathscr M$ guarantees that these local sheaves agree on overlaps, so they glue.  The construction is inverse to $f_*$ by the affine equivalence on each $V$.
\end{proof}

\begin{formalization}
For relative $\Spec$, formalize first on a basis of affine opens closed under suitable refinements by principal opens.  The sheaf condition for $\A$ supplies the cocycle condition for gluing.  The category equivalence in Theorem \ref{thm:affine-qc-equivalence} is best proved by reducing all unit and counit maps to the standard affine module--sheaf equivalence.
\end{formalization}

\section{Vector bundles and locally free sheaves}

\begin{strategy}
Hartshorne's convention associates to a locally free sheaf $\E$ the relative affine scheme
\[
  V(\E)=\Spec_Y\Sym(\E).
\]
With this convention, the sheaf of sections of $V(\E)$ is $\E^\vee$, not $\E$.  This is the main bookkeeping point.  The geometric vector-bundle structure is local: after choosing a basis for $\E$ on an affine open, $\Sym(\E)$ becomes a polynomial algebra, and changes of basis act linearly on the variables.
\end{strategy}

\begin{definition}
A geometric vector bundle of rank $n$ over a scheme $Y$ is a morphism $f:X\to Y$ together with an open cover $\{U_i\}$ of $Y$ and isomorphisms over $U_i$
\[
  \psi_i:f^{-1}(U_i)\xrightarrow{\sim} \mathbb A^n_{U_i}
\]
such that on every affine open $V=\Spec A\subseteq U_i\cap U_j$, the transition automorphism of $\mathbb A^n_V=\Spec A[t_1,\ldots,t_n]$ is induced by a linear change of coordinates
\[
  t_i\longmapsto \sum_j a_{ij}t_j,
\]
with $(a_{ij})\in \operatorname{GL}_n(A)$.
\end{definition}

\begin{proposition}[A locally free sheaf defines a geometric vector bundle]\label{prop:locfree-to-vectorbundle}
Let $\E$ be locally free of rank $n$ on $Y$.  Set
\[
  X=V(\E):=\Spec_Y\Sym(\E).
\]
Then $X\to Y$ is a geometric vector bundle of rank $n$.  Its isomorphism class does not depend on the choice of local bases used to construct the trivializations.
\end{proposition}

\begin{proof}
Let $U=\Spec A\subseteq Y$ be an affine open on which $\E$ is free with basis $e_1,\ldots,e_n$.  Then
\[
  \Sym(\E)(U)\cong \Sym_A(Ae_1\oplus\cdots\oplus Ae_n)
  \cong A[t_1,\ldots,t_n].
\]
Therefore
\[
  f^{-1}(U)=\Spec\Sym(\E)(U)\cong \mathbb A^n_U.
\]

On an overlap, changing the basis of $\E$ by an invertible matrix changes the polynomial generators by the corresponding linear transformation.  Thus the transition maps are linear in the required sense.  Choosing different bases merely composes the local trivializations with such linear automorphisms, so the resulting geometric vector bundle is isomorphic to the original one.
\end{proof}

\begin{proposition}[The sheaf of sections of a vector bundle]\label{prop:sections-sheaf}
Let $f:X\to Y$ be a geometric vector bundle of rank $n$.  The presheaf
\[
  U\longmapsto \{\text{sections } U\to X \text{ over } U\}
\]
is a sheaf of sets, denoted $\mathscr S(X/Y)$, and it carries a natural $\OO_Y$-module structure making it locally free of rank $n$.
\end{proposition}

\begin{proof}
The sheaf property follows because morphisms of schemes glue.  The module structure is local on $Y$.  On a trivializing affine open $U=\Spec A$ with
\[
  f^{-1}(U)\cong \mathbb A^n_U=\Spec A[t_1,\ldots,t_n],
\]
a section $s:U\to \mathbb A^n_U$ corresponds to an $A$-algebra homomorphism
\[
  A[t_1,\ldots,t_n]\longrightarrow A,
\]
hence to the $n$-tuple
\[
  (s_1,\ldots,s_n)=(\theta(t_1),\ldots,\theta(t_n))\in A^n.
\]
Use this identification to define addition and scalar multiplication.  On overlaps, the transition maps are linear, so these local module structures agree and glue.  Locally the section sheaf is $\OO_U^{\oplus n}$, hence it is locally free of rank $n$.
\end{proof}

\begin{proposition}[Sections of $V(\E)$]\label{prop:sections-of-ve}
Let $\E$ be locally free of rank $n$ on $Y$, and let
\[
  X=V(\E)=\Spec_Y\Sym(\E).
\]
Then the sheaf of sections of $X\to Y$ is canonically isomorphic to $\E^\vee$:
\[
  \mathscr S(X/Y)\cong \E^\vee.
\]
\end{proposition}

\begin{proof}
For an affine open $U\subseteq Y$, the ordinary affine scheme--ring correspondence identifies a section $U\to X$ with an $\OO_Y(U)$-algebra homomorphism $\Sym(\E)(U)\to\OO_Y(U)$. These identifications respect restriction to principal opens. Covering an arbitrary open $V\subseteq Y$ by affine opens and gluing therefore identifies sections $V\to X$ over $V$ with $\OO_V$-algebra homomorphisms
\[
  \Sym(\E|_V)\longrightarrow \OO_V.
\]
By the universal property of the symmetric algebra, such homomorphisms are naturally equivalent to $\OO_V$-linear maps
\[
  \E|_V\longrightarrow \OO_V,
\]
that is, to sections of $\E^\vee$ over $V$.  This identification is functorial in $V$ and respects the module structures, so it gives an isomorphism of sheaves of $\OO_Y$-modules.
\end{proof}

\begin{theorem}[Equivalence between locally free sheaves and geometric vector bundles]\label{thm:vectorbundle-correspondence}
There is a bijection between isomorphism classes of locally free sheaves of rank $n$ on $Y$ and isomorphism classes of geometric vector bundles of rank $n$ over $Y$.
\end{theorem}

\begin{proof}
Given a locally free sheaf $\E$, Proposition \ref{prop:locfree-to-vectorbundle} constructs $V(\E)$.  By Proposition \ref{prop:sections-of-ve}, the sheaf of sections of $V(\E)$ is $\E^\vee$.

Conversely, given a geometric vector bundle $X\to Y$, Proposition \ref{prop:sections-sheaf} gives a locally free sheaf $\mathscr S(X/Y)$ of rank $n$.  Applying the construction of Proposition \ref{prop:locfree-to-vectorbundle} to the dual sheaf $\mathscr S(X/Y)^\vee$ recovers $X$.  Indeed, locally on a trivializing affine open, both are $\mathbb A^n_U$, and the transition functions are the same linear transition functions, viewed on coordinate functions.

Thus the two procedures are inverse up to the harmless dual dictated by Hartshorne's convention $V(\E)=\Spec\Sym(\E)$.  Equivalently, one may send a locally free sheaf $\F$ to $V(\F^\vee)$ if one wants the associated sheaf of sections to be literally $\F$.
\end{proof}

\begin{formalization}
The key definitional choice is whether the total space of a vector bundle attached to $\E$ is $\Spec\Sym(\E)$ or $\Spec\Sym(\E^\vee)$.  Hartshorne uses the former in this exercise, so sections identify with $\E^\vee$.  A Lean development should make this convention explicit in the name of the construction.
\end{formalization}

\section*{Lean-port checklist}
The arguments above can be separated into reusable formal lemmas:

\begin{enumerate}[label=\textup{(\arabic*)}]
  \item Nonempty connected noetherian schemes with domain local rings are integral.
  \item For $X=\Proj S$, sections of $\OO_X(n)$ embed as intersections of the chart modules $x_i^nS_{(x_i)}$.
  \item Intersections of normal domains inside a common fraction field are normal.
  \item Veronese subrings of normal $\mathbb N$-graded domains are normal.
  \item On an affine noetherian scheme, quasi-coherent sheaves are directed unions of coherent subsheaves.
  \item The image-definition filtrations for symmetric and exterior powers have the expected associated graded pieces after local splitting.
  \item Relative $\Spec$ glues from the affine spectra $\Spec\A(V)$ along principal-open localizations.
  \item Under Hartshorne's convention $V(\E)=\Spec\Sym(\E)$, the section sheaf is $\E^\vee$.
\end{enumerate}

\end{document}
